% =============================================================================
% Included live section owner -- vacuum offsets and late acceleration
% The section preserves compatibility labels while separating determinant
% algebra, a conditional trace-free completion, ordinary scalar--tensor
% proofs of class, and observed GR-normalised late-source bookkeeping.
% =============================================================================

\section{Vacuum Offsets, Trace-Free Completion, and Late Acceleration}
\label{sec:cc_problem}

\subsection{The problem and the four statements that must be separated}
\label{subsec:cc_statement}

In general relativity coupled to quantum field theory, a constant term in
the matter effective action has the metric dependence
\begin{equation}
 S_\Lambda=-\int d^4x\,\sqrt{-g}\,\rho_\Lambda ,
 \label{eq:cc_standard}
\end{equation}
and therefore sources the metric equations.  The observed late-time source
is conventionally written
\begin{equation}
 \rho_{\Lambda,\mathrm{obs}}
 =3\Omega_{\Lambda0}\,\bar M_{\mathrm{Pl}}^2H_0^2
 \simeq 10^{-47}\,\mathrm{GeV}^4,
 \label{eq:rho_obs}
\end{equation}
whereas unprotected zero-derivative contributions to an effective action are
sensitive to much higher scales.  This is the standard vacuum-energy
naturalness problem~\cite{Weinberg1989,Carroll2001,Polchinski2006}.

Four logically different statements are relevant in ECT and are kept
separate throughout this section.
\begin{enumerate}
\item The orientation kinetic tensor has an exactly orientation-independent
determinant at fixed amplitude and fixed coefficients.  This is an algebraic
result of the ordered-branch ansatz.
\item A separately supplied trace-free/fixed-measure scalar--tensor action is
classically covariant under a constant shift of the scalar potential when an
integration density is shifted oppositely.  This is an exact result inside
that declared completion.
\item The orientation determinant identity alone does not derive the
trace-free metric action, its Bianchi completion, the matter/photon metric,
or the integration density.
\item A regular late-accelerating scalar--tensor proof of class exists, but
its action, state and dimensional calibration are not yet selected by the
frozen ECT postulates.  The observed late source, $H_0$, and the age of our
condensate are therefore matching/state questions until that map is derived.
\end{enumerate}

The standard trace-free algebra is stated below as a conditional completion;
it is not attributed to the orientation identity without the missing
action-level descent.

\subsection{Three source terms and their owners}
\label{subsec:three_objects}

It is useful to distinguish three quantities that are often called vacuum
energy.

\paragraph{Microscopic baseline.}
A field-independent additive constant in the fundamental Euclidean
Lagrangian changes the value of the Euclidean action.  Its effect on
Lorentzian observables depends on the derived continuation, measure and
physical metric action; the absence of an independent Euclidean metric
variation is not by itself a proof that the corresponding Lorentzian source
vanishes.

\paragraph{Zero-derivative quantum term.}
Integrating out fluctuations can generate local terms with field-dependent
prefactors.  Their contribution is governed by the full constrained
Hessian, measure and counterterms, not by the determinant of one kinetic
block alone.

\paragraph{Late-time source.}
The source that enters the physical Friedmann and photon-redshift equations
may be an integration density, a dynamical scalar contribution, another
ordered-branch sector, or a combination of them.  Its law and observable map
must be derived from one complete action; its numerical value for our
condensate may then be fixed as boundary/initial data or inferred from
observations, just as other cosmological state data are.

\subsection{Scope of Weinberg's no-go result}
\label{subsec:weinberg_cc}

Weinberg's adjustment no-go result~\cite{Weinberg1989} is not evaded merely
because the microscopic description is Euclidean or because a Lorentzian
metric is emergent.  Once a local low-energy action, a physical metric and
the adjustment fields are specified, the theorem's hypotheses must be
checked on that completed theory.  The orientation-determinant identity does
not remove those hypotheses and therefore is not a solution of the
cosmological-constant problem.  A separately declared trace-free action can
change the local variational premises and replace a Lagrangian offset by a
global integration datum, but ECT has not yet derived that action, its
quantum measure or the rule selecting the datum.  The two-slope ordinary
completion changes the late dynamics but does not evade the additive-offset
problem at all.  Thus Weinberg's result remains the relevant comparison and
consistency test; whether a complete ECT action lies outside one of its
assumptions is Open rather than established.

\subsection{Exact orientation determinant and its scope}
\label{subsec:main_theorem}

For the fixed-coefficient orientation tensor
\begin{equation}
 K^{AB}=\beta\delta^{AB}-\alpha n^An^B,
 \qquad \delta_{AB}n^An^B=1,
 \label{eq:KAB_cc}
\end{equation}
the matrix determinant lemma gives
\begin{equation}
 \det K^{AB}=\beta^3(\beta-\alpha),
 \label{eq:detK}
\end{equation}
independent of the direction of $n^A$.  If, in addition, this tensor is
identified with a densitized inverse metric through
\begin{equation}
 K^{AB}=C_K\sqrt{-g_{\rm eff}}\,g_{\rm eff}^{AB}
 \quad (C_K=\text{field independent}),
 \label{eq:K_geff_identification}
\end{equation}
then the determinant of that particular representative is fixed,
\begin{equation}
 \sqrt{-g_{\rm eff}}\propto\sqrt{-\det K}
 =\sqrt{\beta^3(\alpha-\beta)},
 \label{eq:sqrt_g_const}
\end{equation}
and tangent orientation variations obey
\begin{equation}
 \delta_n\sqrt{-g_{\rm eff}}=0.
 \label{eq:delta_sqrt_g}
\end{equation}

\paragraph{Exact lemma and non-converse.}
\phantomsection\label{thm:cc_decoupling}% compatibility label
Equations~\eqref{eq:detK}--\eqref{eq:delta_sqrt_g} are Level~A algebra
inside the fixed-$\alpha,\beta$ orientation ansatz.  Their converse is
false: a constant determinant does not determine the metric equations.  For
example,
\begin{equation}
 g_{\mu\nu}(s)=\operatorname{diag}
 (-e^{3s},e^{-s},e^{-s},e^{-s}),\qquad \det g=-1,
 \label{eq:A_decoupling}
\end{equation}
has the same determinant for every $s$ but different lapse, spatial rulers,
redshifts and response.  Moreover, the orientation sphere supplies only
three tangent variations, while a trace-free symmetric metric contains nine
components before diffeomorphism constraints.  Amplitude dependence,
orientation--amplitude mixing and field-dependent normalization can also
change the physical measure while leaving Eq.~\eqref{eq:detK} intact.

Consequently the exact orientation lemma is not an unconditional
vacuum-offset decoupling theorem.  Deriving the physical constrained metric
action, its constraint algebra, its Bianchi identity and its universal
matter/photon coupling remains an Open action-level task.

\subsection{Conditional trace-free/fixed-measure completion}
\label{subsec:unimodular_cc}

The standard unimodular mechanism can nevertheless be stated exactly once a
specific completion is supplied.  Consider
\begin{equation}
 S_{\rm TF}=\int d^4x\left\{\sqrt{-g}\left[
 {1\over2}f(q)R-{1\over2}K(q)(\partial q)^2-V(q)
 +\mathcal L_m\right]
 +\lambda(x)\bigl(\sqrt{-g}-\epsilon_0\bigr)\right\}.
 \label{eq:full_quadratic_cc}
\end{equation}
Variation with respect to $\lambda$ fixes the measure.  Metric variation,
the scalar equation, matter conservation and the Bianchi identity imply
$\lambda=-C_I$ with $\partial_\mu C_I=0$, and reconstruct
\begin{align}
 fG_{\mu\nu}={}&T^m_{\mu\nu}
 +K\left(\partial_\mu q\partial_\nu q
 -{1\over2}g_{\mu\nu}(\partial q)^2\right)
 -[V(q)+C_I]g_{\mu\nu}\nonumber\\
 &+\nabla_\mu\nabla_\nu f-g_{\mu\nu}\Box f.
 \label{eq:H1_cc}
\end{align}
The scalar equation contains $V_{,q}$ but not $C_I$.  Therefore
\begin{equation}
 V(q)\longrightarrow V(q)+C_0,
 \qquad C_I\longrightarrow C_I-C_0
 \label{eq:B_decoupling}
\end{equation}
leaves the complete classical equations invariant.  This is the precise
conditional vacuum-offset result: only the combination $V+C_I$ enters the
metric equation, while the split is conventional.

In a constant-scalar vacuum,
\begin{equation}
 R={4[V(q)+C_I]\over f(q)},\qquad
 V_{,q}-2{f_{,q}\over f}[V(q)+C_I]=0.
 \label{eq:H1_broken_correction}
\end{equation}
For $f=f_*e^{aq}$ and $V=V_1e^{bq}$, the root is
\begin{equation}
 e^{bq_0}={2aC_I\over(b-2a)V_1}.
 \label{eq:cc_tracefree_root}
\end{equation}
It exists as a real positive exponential precisely when
$aC_I/[(b-2a)V_1]>0$ (with nonzero denominator).  Choices such as
$C_I>0$, $V_1>0$, $a>0$ and $b>2a$ are one branch, not the general
existence condition.

\paragraph{Interpretation of $C_I$.}
$C_I$ is a global integration/state datum, not a value determined by the
local equations.  This does not require a fundamental law to calculate the
particular datum of our condensate: it may be selected by boundary/initial
conditions and measured.  What ECT still owes is the derivation of
Eq.~\eqref{eq:full_quadratic_cc} (or another complete shift-covariant action)
from the $\Phi$ medium and the map by which $C_I$ enters physical clocks,
matter and photons.  Until then, the trace-free result is Level~A within the
named completion and Level~B/Open as an ECT descent.

\subsection{Why the fixed determinant does not supply the completion}
\label{subsec:Hlambda_cc}

The earlier spectral condition $H_\Lambda$ can be represented schematically
by a quadratic fluctuation operator
\begin{equation}
 S^{(2)}={1\over2}\int d^4x_E\left[
 \mathcal K^{AB}(\phi,n,\ldots)\partial_A\chi\partial_B\chi
 +\mathcal M^2(\phi,n,\ldots)\chi^2\right].
 \label{eq:full_quadratic_cc_compat}
\end{equation}
Even if
\begin{equation}
 \mathcal K^{AB}=C\,K^{AB}(n),\qquad C=\text{constant},
 \label{eq:H1_cc_compat}
\end{equation}
the determinant controls only that Gaussian block.  It does not establish
the physical constrained metric action or a radiative
non-renormalization theorem.  If instead
$\mathcal K^{AB}=Z(q)K^{AB}$, then
\begin{equation}
 \delta_q\Gamma_{\rm vac}^{(0)}
 \propto\Lambda_{\rm UV}^4Z(q)Z_{,q}(q),
 \label{eq:H1_broken_correction_compat}
\end{equation}
and a scalar naturalness problem is explicit.  Conditions on the full
reduced Hessian, measure, constraints and counterterms are therefore needed
even after the classical trace-free action is derived.

\subsection{Ordinary two-slope completion: regular dynamics, not shift
invariance}
\label{subsec:frw_cc}

A distinct constructive result uses the ordinary Jordan action
\begin{equation}
 S_J=\int d^4x\sqrt{-g}\left[
 {1\over2}f(q)R-{1\over2}K(q)(\partial q)^2-V(q)\right]+S_m[g,\psi],
 \quad
 f=f_*e^{aq},\quad K=\kappa f,\quad C_I=0,
 \label{eq:frw_lagrangian_cc}
\end{equation}
with
\begin{equation}
 V(q)=V_-e^{cq}+V_+e^{bq},\qquad
 0<c<2a<b,\qquad V_\pm>0.
 \label{eq:cc_two_slope_action}
\end{equation}
Its Einstein-frame potential
$\mathcal U_E\propto V_-e^{(c-2a)q}+V_+e^{(b-2a)q}$ is coercive and
has one strict minimum,
\begin{equation}
 e^{(b-c)q_{\rm vac}}
 ={(2a-c)V_-\over(b-2a)V_+},\qquad
 \mathcal E_{,q}(q_{\rm vac})
 =(2a-c)(b-c)V_-e^{cq_{\rm vac}}>0.
 \label{eq:cc_two_slope_root}
\end{equation}
Provided $f>0$, $\kappa+3a^2/2>0$, and the solution is on the expanding
finite-energy branch, this supplies a
regular late-de-Sitter proof of class and, for a frozen
near-GR benchmark, a connected radiation--matter--de-Sitter numerical
history.  It does not make the ordinary action invariant under
$V\to V+C_0$: an added constant still gravitates and changes the scalar
balance.  The two-slope mechanism and the trace-free offset mechanism are
therefore complementary, not interchangeable.

Here $\kappa+3a^2/2>0$ is equivalently the positivity condition for the
Einstein-frame scalar kinetic coefficient in the displayed normalisation.

For a homogeneous scalar the bare Jordan-bookkeeping scalar pieces are
\begin{align}
 \rho_q&={1\over2}K\dot q^2+V(q),
 \label{eq:rho_q_cc}\\
 p_q&={1\over2}K\dot q^2-V(q).
 \label{eq:p_q_cc}
\end{align}
They are neither separately conserved in general nor the full effective
dark-energy density and pressure: non-minimal-coupling terms proportional to
$\dot f$, $\ddot f$ and the chosen metric rearrangement also enter the
Friedmann bookkeeping.  The ratio $p_q/\rho_q\ge-1$ follows for these bare
pieces only when $K\ge0$ and $\rho_q>0$; it is not a prediction of the
physical ECT dark-energy equation of state.

\paragraph{Current observational context and a conditional falsifier.}
DESI DR2 BAO are well described by flat $\Lambda$CDM in isolation.  When
combined with CMB data and, dataset-dependently, supernova compilations,
$w_0$--$w_a$ fits prefer evolving dark energy, with the quoted significance
depending on the combination.  The extended DESI analysis finds a common
low-redshift trend and a preference for phantom crossing, while non-crossing
alternatives are disfavoured but not excluded
~\cite{DESIDR2_2025,Lodha2025}.  This is observational context, not an ECT
prediction or a discovery claim.  In the positive-kinetic two-slope Jordan
scalar the bare ratio
$w_q=(K\dot q^2/2-V)/(K\dot q^2/2+V)$ obeys $w_q\ge-1$ when $K\ge0$ and
$\rho_q>0$, but the nonminimal-$f$ terms mean that this bare ratio is not the
effective $w(z)$ fitted by cosmological data.  A complete action-owned ECT
history robustly constrained never to cross $-1$ would therefore be
disfavoured by that combined-data trend.  The present candidate has not
supplied the physical photon/perturbation likelihood and is untested by this
falsifier.

\subsection{Observed late source, $H_0$, and age}
\label{subsec:IR_estimate_cc}

After choosing the present physical metric and rewriting its measured
background in the GR-normalised form with fixed reference
$\bar M_{\rm Pl}$, define
\begin{equation}
 \rho_{\rm late,0}^{(\rm GR\,ref)}
 \equiv3\Omega_{\rm late,0}^{(\rm GR\,ref)}
 \bar M_{\rm Pl}^2H_0^2.
 \label{eq:IR_estimate_cc}
\end{equation}
This is an observational convention/definition, not the generic
scalar--tensor Friedmann equation and not a dimensional derivation.  In the
supplied scalar--tensor action the present constraint also contains $F_0$ and
$-3H_0\dot F_0$ (and any separately supplied orientation stress); how those
terms are assigned to an effective late component must be stated.  Writing
$\rho_{\rm late,0}^{(\rm GR\,ref)}=
c_\Lambda^{\rm GR\,ref}\bar M_{\rm Pl}^2H_0^2$ therefore defines
$c_\Lambda^{\rm GR\,ref}=3\Omega_{\rm late,0}^{(\rm GR\,ref)}$
only in this GR-reference bookkeeping
after $H_0$, the physical metric and Planck calibration are fixed.  It is not
an ECT determination or evidence that a microscopic vacuum baseline has
decoupled.

As a measured-input dimensional check,
$\bar M_{\rm Pl}=2.435\times10^{18}\,\mathrm{GeV}$ and
$H_0\simeq1.44\times10^{-42}\,\mathrm{GeV}$ give
$\bar M_{\rm Pl}^2H_0^2\simeq1.2\times10^{-47}\,\mathrm{GeV}^4$, of the
same dimensional order as the observed late source.  This arithmetic uses
measured inputs and is not a derivation of either the scale or its
coefficient from ECT.

The dimensional alternatives also diagnose what has and has not been
explained.  A density proportional to $\bar M_{\rm Pl}H_0^3$ is smaller than
$\bar M_{\rm Pl}^2H_0^2$ by $H_0/\bar M_{\rm Pl}$, while $H_0^4$ is smaller
by $(H_0/\bar M_{\rm Pl})^2$; neither can replace the observed scale without
an additional large dimensionless factor or new scale.  The
Cohen--Kaplan--Nelson UV/IR consistency estimate
$\rho\lesssim \bar M_{\rm Pl}^2/L_{\rm IR}^2$
~\cite{CohenKaplanNelson1999} reproduces the same dimensional order only
after $L_{\rm IR}\sim H_0^{-1}$ and near saturation are supplied.  It is a
useful external comparison, not an ECT derivation.

The two-slope proof-of-class yields, for one explicitly calibrated
near-GR benchmark,
\begin{equation}
 \tau_{2s}\equiv H_{2s}(0)\Delta t=0.96368677864.
 \label{eq:cc_conditional_age}
\end{equation}
Using measured values of $H_0$ as two disclosed alternatives for the single
dimensional scale calibration of this already fixed orbit gives
$13.98050$~Gyr for $67.4\,\mathrm{km\,s^{-1}\,Mpc^{-1}}$ and
$12.90095$~Gyr for $73.04\,\mathrm{km\,s^{-1}\,Mpc^{-1}}$.  These are
conditional outputs of one calibrated action and state, not unique ECT
predictions.  Different regular completions or values of $C_I$ give
different histories.  Conversely, leaving $C_I$ as a measured state datum
is logically acceptable once its action-level and observable maps are
derived.

\subsection{Infrared carrier remains unidentified}
\label{subsec:goldstone_cc}

The current action derives no physical soft-mode operator, mass, state or
abundance that carries the observed late source.  A pseudo-Goldstone
interpretation would require all of those owners and a cosmological
forward model.  Writing $H_{0,E}:=\hbar H_0^{(\mathrm{freq})}$ so that the
physical frequency and its natural-unit energy representation are explicit,
the separately named C/Open scale hypothesis
$m_{\chi,E}=\zeta_HH_{0,E}$ may be used only for conditional orientation;
it is not selected by the action and supplies no late-source identification.
No Planck-scale field excursion, numerical infrared-density scaling,
abundance or cosmological-role inference is retained.

\subsection{Comparison and status}
\label{subsec:comparison_cc}

The corrected result should also be compared with established external
routes rather than presented in isolation:

{\small
\setlength{\tabcolsep}{4.5pt}
\setlength{\LTpre}{0pt}
\setlength{\LTpost}{0pt}
\renewcommand{\arraystretch}{0.98}
\begin{longtable}{p{3.0cm}p{4.0cm}p{3.5cm}p{\dimexpr3.6cm+12pt\relax}}
\caption{Compact comparison with external vacuum-energy approaches.}
\label{tab:cc_external_approaches}\\
\toprule
\textbf{Approach} & \textbf{What it changes} & \textbf{What remains input}
& \textbf{Relation to the present ECT programme}\\
\midrule
GR plus QFT & A constant term sources the metric & Renormalized residual
vacuum density & Baseline naturalness problem~\cite{Weinberg1989}\\
Unimodular/trace-free gravity & Local trace equation is replaced; a global
integration constant appears & Boundary/state value and quantum completion
& Mathematical template for the conditional $C_I$ branch
~\cite{HenneauxTeitelboim1989}\\
Sequestering/global constraints & Global variables cancel selected vacuum
contributions & Complete global action, measure and boundary prescription
& Possible external completion class, not derived here\\
Dynamical scalar/quintessence & A rolling field supplies late acceleration
& Potential, initial state, matter metric and radiative stability & The
ordinary two-slope branch is a conditional member of this broad class\\
Anthropic/landscape selection & Restricts attention to vacuum values
compatible with observers or structure & Ensemble, measure and selection
rule & External comparison only; not an ECT derivation
~\cite{Weinberg1987}\\
CKN UV/IR bound & Relates an infrared size to an allowed energy density
& Choice of $L_{\rm IR}$ and degree of saturation & Dimensional comparison
only~\cite{CohenKaplanNelson1999}\\
\bottomrule
\end{longtable}}

{\small
\begin{longtable}{p{3.1cm}p{3.7cm}p{3.6cm}p{3.7cm}}
\caption{Vacuum-offset and late-source status.  ``Conditional'' always means
inside the explicitly named completion, not a result of the orientation
determinant alone.}
\label{tab:cc_approaches}\\
\toprule
\textbf{Structure} & \textbf{Exact result} &
\textbf{Missing owner} & \textbf{Status}\\
\midrule
Orientation tensor & $\det K$ is direction independent at fixed
$\alpha,\beta$ & Physical metric action and measure & Level~A algebra;
descent Open\\
Explicit trace-free action & $V\!+\!C_I$ is shift covariant; $C_I$ is a
global datum & Derivation from the $\Phi$ medium and physical clock/photon
map & Level~A conditional; ECT map Level~B/Open\\
Ordinary two-slope action & Coercive potential, stable vacuum and regular
background proof of class & Selection of action, parameters and state;
not shift invariant & Level~A conditional / Level~C benchmark\\
Observed late source & Friedmann normalization after $H_0$, metric and state
are fixed & One complete ECT action and state-to-observable map & Open as a
unique ECT cosmology; measurable state datum allowed\\
\bottomrule
\end{longtable}}

\subsection{Binding hierarchy and open tests}
\label{subsec:levels_cc}

\paragraph{Established.}
The orientation determinant identity is exact in its stated ansatz.  The
shift covariance of $V+C_I$ is exact in the separately supplied trace-free
completion.  The coercive two-slope potential and its stable vacuum are
exact in the separately supplied ordinary scalar--tensor completion.

\paragraph{Not established.}
The fixed orientation determinant does not derive the trace-free action;
the microscopic baseline and loop contribution have not been shown to be
absent from the complete physical ECT metric equations; the physical
matter/photon metric is not identified by this section; and the observed
late-time source is not a unique output of the frozen postulates.

\subsection{Open programme}
\label{subsec:cc_open}

\begin{description}
\item[O1.] Derive, or refute, the trace-free/fixed-measure action and its
Bianchi completion from constrained $\Phi$-medium variables.  The derivation
must cover amplitude as well as orientation variations.
\item[O2.] Derive the universal physical metric and show how the integration
density or dynamical late source enters clocks, redshift, lensing and PPN
observables.
\item[O3.] Compute the full reduced Hessian, measure and counterterms before
claiming radiative protection of a zero-derivative term.
\item[O4.] Decide whether the two-slope completion, the trace-free completion,
or a coupled version follows from the ordered $\Phi$ phase.  If $C_I$ is a
state datum, identify its boundary/state-selection rule; its particular
value need not be a universal constant predicted by local postulates.
\item[O5.] Fit one action-owned background and perturbation system jointly to
expansion, growth, CMB/lensing and local-gravity data.  Conditional age
estimates remain valid benchmarks but are not promoted by the fit alone.
\end{description}

\subsection{Summary}
\label{subsec:cc_summary}

The robust result carried by the ordered orientation tensor is the fixed
determinant of one fixed-coefficient kinetic block.  It is not, by itself, a
vacuum-energy decoupling theorem.  An explicit trace-free/fixed-measure
scalar--tensor completion does have an exact classical shift covariance:
$V\to V+C_0$ is absorbed by $C_I\to C_I-C_0$.  The integration density is a
legitimate global state/boundary datum, but ECT has not yet derived the
completion or its observable map.  Separately, an ordinary two-slope action
provides a stable and globally regular late-accelerating proof of class, but
it is not invariant under an added constant and its numerical history is a
calibrated conditional result.  Hence the vacuum-offset route is
constructive, while the physical late source, $H_0$ and the age of our
condensate remain to be identified by one complete action, state and metric
map.

% =============================================================================
% END INCLUDED CC LATE-SOURCE OWNER; R177 SUCCESSOR CANDIDATE
% =============================================================================
